\section{Implications for dark matter direct detection}
\label{sec: direct detection}

In this section we discuss the impact of the LMC on the interpretation of the results of DM direct detection experiments. In particular, in sections~\ref{sec: nuclear} and \ref{sec: electron} we consider the DM interaction with a target nucleus or electron, respectively, and study how the exclusion limits set by different direct detection experiments in the plane of the DM mass and scattering cross section change due to the presence of the LMC for a given experimental setup.

We simulate the signals in three different idealized direct detection experiments, which are inspired by near future detectors that would search for nuclear  or electron recoils due to the scattering with a DM particle. In order to find the constraints in the DM scattering cross section and mass plane, we employ the Poisson likelihood method  implemented in the {\rm DDCalc}~\cite{ddcalc} and {\rm QEDark}~\cite{QE, essig} software packages for nuclear and electron recoils, respectively. Taking the properties of the experiments and the local DM distribution as inputs, these packages provide the exclusion limits at a desired confidence level. To perform the direct detection analysis, we directly use the local DM velocity distributions extracted from the simulations.




\subsection{Dark matter - nucleus scattering}
\label{sec: nuclear}



In the case of DM-nucleus scattering, we consider a DM particle of mass $m_\chi$ scattering with a target nucleus of mass $m_T$ in an underground detector, and depositing the nuclear recoil energy, $E_R$. The differential events rate is given by
\begin{equation}
\frac{dR}{dE_R} = \frac{\rho_\chi}{m_\chi}\frac{1}{m_T} \int_{v>v_{\rm min}} d^3 v ~\frac{d\sigma_T}{d E_R}~v~\tilde{f}_{\rm det}({\bf v}, t)\; ,
\label{eq:difrate}
\end{equation}
where $\sigma_T$ is the DM-nucleus scattering cross section. 
Assuming elastic scattering, the minimum  speed required for a DM particle to deposit a recoil energy $E_R$ to the detector is given by
\begin{equation}
v_{\rm min}(E_R) = \sqrt{\frac{m_T E_R}{2 \mu^2_{\chi T}}}\; ,
\label{eq: vmin}
\end{equation}
where $\mu_{\chi T}$ is the reduced mass of the DM and nucleus.

For spin-independent interactions, the differential cross section is given by
\begin{equation}
  \frac{d\sigma_T}{dE_R}=\frac{m_T A^2 \sigma^{\rm SI}_{\chi N}}{2\mu_{\chi N}^2v^2}\,F^2(E_R)\; ,
\end{equation}
where $A$ is the atomic mass number of the target nucleus, $\sigma^{\rm SI}_{\chi N}$ is the spin-independent DM-nucleon scattering cross section at zero momentum transfer, $\mu_{\chi N}$ is the reduced mass of the DM and nucleon, and $F(E_R)$ is the spin-independent nuclear form factor for which we use the Helm form factor~\cite{Helm:1956zz}.

Hence, the differential event rate can be written in terms of the halo integral (eq.~\eqref{eq: halo integral}) as
\begin{equation}
    \frac{dR}{dE_R} = \frac{\rho_\chi A^2 \sigma^{\rm SI}_{\chi N}}{2 m_\chi \mu_{\chi N}^2}\,F^2(E_R)\,\eta (v_{\rm min}, t)\, .
\end{equation}



We consider two idealized direct detection experiments, one with a xenon target nucleus and the other with germanium. These detectors are based on the sensitivity of the  LZ~\cite{LZ:2022lsv, lz2015} and SuperCDMS~\cite{supercdms} direct detection experiments in the near future. Noble liquid detectors, such as LZ which has recently published its first results~\cite{LZ:2022lsv}, can reach large exposures and are sensitive to large DM masses and lower cross sections. On the other hand, cryogenic solid state detectors such as SuperCDMS are sensitive sub-GeV DM masses. Considered together, these two types of experiments probe a large range of DM masses and scattering cross sections.





For the xenon based experiment, we consider an energy range of $[2 -50]$~keV, an energy resolution of $\sigma_E=0.065 E_R+0.24\,{\rm keV}\sqrt{E_R/{\rm keV}}$~\cite{Bertone:2017adx}, and an exposure of $5.6\times 10^6$~kg~days with a maximum efficiency of 50\% as given in ref.~\cite{lz2015}. 
The exposure we consider for this experiment is expected to be achieved by LZ after five years of operation~\cite{lz2015}.

For the germanium based experiment, we consider two crystal target designs with different energy thresholds. The low energy threshold design is based on the projected high-voltage (HV) detector of the SuperCDMS SNOLAB experiment~\cite{supercdms}. We implement an energy  range of $[40 - 300]$~eV, with a constant signal efficiency of 85\%, a flat background level of 10 keV$^{-1}$~kg$^{-1}$~days$^{-1}$, and an exposure of $1.6\times10^4$~kg~days~\cite{supercdms, Kahlhoefer:2017ddj}. The high energy threshold design has similar features as the iZIP detector of the same experiment with a total exposure of $2.04\times10^4$~kg~days, an energy range of $[3 - 30]$~keV, 1 expected background event, and a flat efficiency of 75\%. The exposures considered are expected to be achieved by SuperCDMS after five years of operation~\cite{supercdms}.

\begin{figure}[t]
\centering
    \includegraphics[width=\linewidth]{graphics/Xe-grids.pdf}
    \caption{Top panels: exclusion limits at 90\% CL for a future xenon based experiment in the spin-independent DM-nucleon cross section and DM mass plane for four snapshots in halo 13: the isolated MW analogue (Iso.), the LMC's pericenter approach (Peri.),  the present day MW-LMC analogue (Pres.), and the future MW-LMC analogue (Fut.). For each snapshot, the solid/dashed lines and the shaded bands correspond to the exclusion limits obtained from the mean and the $1\sigma$ uncertainty band of the halo integrals, respectively. For the pericenter, present day, and future snapshots, the exclusion limits are presented in the Solar region for the best fit Sun's position for MW+LMC (solid coloured curves) and MW-only (dashed coloured curves) DM populations. For the isolated MW snapshot, the exclusion limit is shown for the DM population of the  MW  (solid black curve)  for a Solar region defined as a spherical shell with radii between 6 and 10 kpc from the Galactic center. The blue curve correspond to the exclusion limit for the SHM Maxwellian. The local DM density is set to $\rho_{\chi} = 0.3$~GeV/cm$^3$. Bottom panels: the ratios of the exclusion limits for the MW-only and the MW+LMC DM populations for the pericenter, present day, and future snapshots. The left panels show the limits and ratios for a large range of DM masses,
while the right panels zoom onto the low DM mass region.}
\label{fig:LZ_bestfit_merged}
\end{figure}

The top panels of figures~\ref{fig:LZ_bestfit_merged} and \ref{fig:supercdms_bestfit}  show the exclusion limits at the 90\% CL in the plane of DM mass and spin-independent cross section set by the future xenon and germanium experiments using the local DM velocity distribution at the four representative  snapshots in halo 13, respectively. These snapshots are: the isolated MW analogue (Iso.), the LMC’s pericenter approach (Peri.),  the present day MW-LMC analogue (Pres.), and the future MW-LMC analogue (Fut.).  The mean and the shaded band in the exclusion limits are obtained from the mean and $1\sigma$ uncertainty band of the halo integrals shown in figure~\ref{fig: halo int response}, respectively. 
The exclusion limit for the isolated MW analogue is shown as the solid black curve, 
while the exclusion limits for the three other snapshots are shown as solid coloured curves for the MW+LMC distribution and dashed coloured curves for the MW-only distribution. For comparison, the exclusion limit assuming the SHM Maxwellian velocity distribution with peak speed of 220~km/s and truncated at the escape speed of 544~km/s from the Galaxy is shown as the solid blue curve. To distinguish the effect of the local DM velocity distribution, the local DM density is set to $\rho_\chi=0.3$~GeV/cm$^3$ in all cases\footnote{Since $\rho_\chi$ is a normalization in the event rate, changing its value results in the same upward or downward shift in all the exclusion limits.}, as is commonly adopted in the SHM.  The bottom panels of the figures show the ratios of the exclusion limits of the MW-only to the MW+LMC 
distribution  for  the pericenter, present day, and future snapshots. The left  panels show the limits and ratios for a large range of DM masses, while the  right panels zoom onto the low DM mass region to better visualize the differences at low masses. 



\begin{figure}[t]
\centering
\includegraphics[width=\textwidth]{graphics/Ge-grids.pdf}
\caption{Same as figure~\ref{fig:LZ_bestfit_merged}, but for a future germanium based experiment.}
\label{fig:supercdms_bestfit}
\end{figure}


The trends in figures~\ref{fig:LZ_bestfit_merged} and \ref{fig:supercdms_bestfit} are similar to those seen in figure~\ref{fig: halo int response} for the halo integrals of the different snapshots. In particular, the differences in the high speed tail of the halo integrals lead to variations in the exclusion limits at low DM masses, where the experiments are most sensitive to high values of $v_{\rm min}$. The isolated MW snapshot has the weakest exclusion limit at low DM masses and follows closely the SHM exclusion limit, while the DM distribution of the MW+LMC  at the LMC's pericenter approach leads to the strongest exclusion limit. As it can be seen from figure~\ref{fig:LZ_bestfit_merged}, for the xenon based experiment the exclusion limit for the MW+LMC distribution at the present day snapshot is lower than the isolated MW exclusion limit by an order of magnitude at $m_\chi \sim 8$~GeV, by more than three orders of magnitude  at $m_\chi \sim 6$~GeV, and by more than five orders of magnitude  at $m_\chi \sim 5$~GeV. Moreover at fixed cross sections, the exclusion limit for the MW+LMC distribution at the present day snapshot shows a shift of a few GeV towards smaller DM masses compared to the isolated MW for masses below $\mathcal{O}(10~{\rm GeV})$. Figure~\ref{fig:supercdms_bestfit} shows that for the germanium based experiment an order of magnitude  of vertical shift  occurs at $m_\chi \sim 0.5$~GeV between the exclusion limits of the MW+LMC distribution at the present day snapshot and the isolated MW, while the vertical shift is more than three orders of magnitude at  $m_\chi \sim 0.4$~GeV. Furthermore, at fixed cross sections and for DM masses below $\mathcal{O}(1~{\rm GeV})$, 
a horizontal shift of a few hundreds of MeV happens towards smaller DM masses. Hence, we can see from figures~\ref{fig:LZ_bestfit_merged} and \ref{fig:supercdms_bestfit} that the LMC extends the parameter space probed by direct detection experiments towards smaller DM masses.

Our results agree with those of ref.~\cite{Besla:2019xbx}, which also found that the presence of the LMC causes direct detection limits to shift to lower cross sections and lower DM masses, extending the sensitivity of those experiments. Hence, our results confirm that the findings of ref.~\cite{Besla:2019xbx} hold even in a fully cosmological setting. 



\subsection{Dark matter - electron scattering}
\label{sec: electron}

In the case of DM-electron scattering, the differential event rate in a crystal target is given by~\cite{Essig:2015cda}
\begin{equation}
    \label{eq:diff_rate}
  \frac{dR}{ d\ln E_{e}  }  =  N_{\rm cell}\frac{\rho_\chi}{m_\chi}\frac{\overline\sigma_e \alpha m_e^2}{\mu_{\chi e}^2} \int d\ln q\, \frac{E_e}{q}\left[ |F_{\rm DM}(q)|^2 \,
              |f^{\rm crystal}(E_e, q)|^2 \,
                   \eta(v_{\rm min}(q, E_e)) \right]\; ,
\end{equation}
where $E_e$ is the energy deposited to the electron, $q$ is the  momentum transfer between the DM and the electron, $N_{\rm cell}$ is the number of unit cells per mass in the crystal target, $\overline\sigma_e$ is the DM-electron reference scattering cross section which parameterizes the strength of the interaction, $\alpha \simeq 1/137$ is the fine structure constant, $m_e$ is the mass of the electron, and $\mu_{\chi e}$ is the reduced DM-electron mass. The dimensionless crystal form factor, $f^{\rm crystal}$, encodes the dependence of the rate on the electronic structure of the target material. 

The DM form factor, $F_{\rm DM}$, gives the momentum dependence of the interaction. It can be shown that $F_{\rm DM}(q)=1$ for a contact interaction via a heavy mediator, $F_{\rm DM}(q)=(\alpha m_e/q)$ for an electron dipole moment coupling, and $F_{\rm DM}=(\alpha m_e/q)^2$ for a long-range interaction induced by the exchange of an ultralight or massless mediator~\cite{Essig:2015cda}.

Lastly, the minimum speed required for the DM particle in order for the electron to gain an energy $E_e$ with momentum transfer $q$ is given by
\begin{equation}
v_{\rm min}(E_e,q)=\frac{E_e}{q}+\frac{q}{2m_\chi}\, .
\label{eq: vmin-e}
\end{equation}



We consider a future silicon CCD experiment, based on the sensitivity of the next generation kg-sized DAMIC-M~\cite{Lee:2020abc, Castello-Mor:2020jhd, DAMIC-M:2023gxo} experiment. Direct detection experiments searching for DM-electron interactions provide a new avenue to probe MeV DM masses, due to the small mass of the electron. Semiconductors, in particular, have a very low ionization threshold of $\sim 1$~eV, and can be sensitive to single electron-hole pairs. We consider a silicon based detector with  an exposure of 1~kg~year and assuming zero background events, with an ionization threshold of 1 electron-hole pair. 

The top panels of figure~\ref{fig:Si_electron_merged} show the exclusion limits at the 95\% CL in the plane of DM mass and DM-electron cross section for the future silicon based experiment, using the local DM velocity distribution at the isolated MW (black), pericenter (green), present day (orange), and future (magenta) snapshots of halo 13. The exclusion limits for the three latter snapshots are shown as
solid coloured curves for the MW+LMC distribution and dashed coloured curves for the MW-only distribution.  The mean and the
shaded band in the exclusion limits are obtained from the mean and $1\sigma$ uncertainty band of
the halo integrals shown in figure~\ref{fig: halo int response}, respectively. The SHM exclusion limit is shown as the solid blue curve. As for the DM-nucleus scattering, the local
DM density is set to $\rho_\chi = 0.3$~GeV/cm$^3$. The bottom panels show the ratio of the exclusion limits of the MW-only  to the MW+LMC distributions for the pericenter, present day, and future snapshots. The left, middle, and right panels show the results for three different DM  form factors, $F_{\rm{DM}}=1$, $F_{\rm{DM}} \propto q^{-1}$, and $F_{\rm{DM}} \propto q^{-2}$, respectively.

The general implications of the LMC for the exclusion limits on the DM-electron scattering cross section are similar to the DM-nucleus scattering, although the effect is smaller in the former case. As it can be seen from figure~\ref{fig:Si_electron_merged}, for all three choices of the DM form factor, the exclusion limits of the MW+LMC distribution at the LMC's pericenter approach and the present day MW-LMC show a shift towards smaller DM masses and lower DM-electron cross sections compared to the isolated MW. As expected, the shift becomes larger for smaller DM masses, where the experiment probes larger $v_{\rm min}$. In particular, the exclusion limit for the MW+LMC distribution at the present day snapshot is lower than the isolated MW exclusion limit by up to a factor $\sim 4$ at $m_\chi \sim 1$~MeV, and by up to a factor $\sim 50$ at $m_\chi \sim 0.6$~MeV. For DM masses below a few MeV and at fixed cross sections, the exclusion limit is shifted by a fraction of MeV towards smaller masses for all three choices of the DM form factor.






\begin{figure}[t]
  \centering
  \includegraphics[width=\textwidth]{graphics/Elec-grids-new.pdf}
  \caption{Top panels: exclusion limits at 95\% CL for a future silicon based experiment in the DM-electron cross section and DM mass plane for four snapshots in halo 13: the isolated MW analogue (Iso.), the LMC's pericenter approach (Peri.),  the present day MW-LMC analogue (Pres.), and the future MW-LMC (Fut.). The DM form factor is assumed to be $F_{\rm{DM}}=1$ (left panel), $F_{\rm{DM}} \propto q^{-1}$ (middle panel), and $F_{\rm{DM}} \propto q^{-2}$ (right panel). The blue curve corresponds to the SHM exclusion limit. The local DM density is set to $\rho_{\chi} = 0.3$~GeV/cm$^3$. Bottom panels: the ratios of 
  the exclusion limits for the MW-only and the MW+LMC  DM populations for the pericenter, present day, and future snapshots. The description of the different coloured curves are the same as in figure~\ref{fig:LZ_bestfit_merged}.
  }
  \label{fig:Si_electron_merged}
\end{figure}



